% Drawing 146/04-10 -- folder 146, batch 04, archive page 67 (PDF page 68).
% Transcription: Plusieurs configurations de courbes stables : une droite portant deux grappes de composantes, marquée $5$ ; deux droites sécantes dont l'une porte trois points et un point marqué $5^*$, $i$, avec la légende $1 \leqslant i \leqslant 4$ ; une droite portant trois composantes verticales ; trois droites dont l'une porte un point, marquée $\mathbb{P}$.
% Kind: configurations of stable curves (schematic). Method 2 (schematic redraw from the facsimile). Status: ok.
% Notes: Four configurations, drawn as on the page: a line with two clusters of components (« 5 »); two crossing lines, one with three points, the other with a point « 5*, i » and the legend 1 <= i <= 4; a line with three vertical components (two of them forked); three lines, one carrying a point, marked P. Not encoded as dual graphs: reading these as Mumford-Deligne graphs is an interpretation the page does not spell out. The crossed pair of pencil lines higher on the page (x 0.58-0.68, y 0.24-0.37) is not part of this mark and is not described by the transcription.
% Coordinates: 1 unit = 1 mm = 1/100 of the crop width (bbox in page fractions: [0.2, 0.53, 0.82, 0.725]).
% Redrawn by Claude (claude-opus-5-5), issue #73 pilot, 2026-11-07; not checked by a human; not compiled (no TeX in the session).
% Rights: reproduces part of the work -- behind RENDER_TRANSCRIPTS, never published by itself.
\documentclass[tikz,border=3pt]{standalone}
\usepackage{amsmath,amssymb}
\usetikzlibrary{patterns,arrows.meta}
\tikzset{pencil/.style={black!45}, struck/.style={black!35}, mark/.style={circle,draw,inner sep=1.3pt}, pt/.style={circle,fill,inner sep=0.9pt}, opt/.style={circle,draw,fill=white,inner sep=1.1pt}}
\begin{document}
\begin{tikzpicture}[x=1mm,y=1mm,line cap=round,line join=round,font=\scriptsize]
\draw (14.5,-9) -- (25,-8.5);
\draw (14.5,-9) -- (13,-3);
\draw (14.5,-9) -- (16,-2.5);
\draw (25,-8.5) -- (24,-3.5);
\draw (25,-8.5) -- (27,-3);
\draw (25,-8.5) -- (28.5,-3.5);
\node[inner sep=1pt] at (11,-3.5) {\scriptsize $5$};
\draw (18,-12) -- (31,-27);
\draw (20.5,-36.5) -- (28.5,-13.5);
\node[pt] at (21.5,-16) {};
\node[pt] at (23,-29.5) {};
\node[pt] at (21.5,-32.5) {};
\node[pt] at (21,-35) {};
\node[inner sep=1pt] at (17.5,-14) {\scriptsize $5^*$};
\node[inner sep=1pt] at (20,-18.5) {\tiny $i$};
\node[inner sep=1pt] at (12,-24) {\scriptsize $1\leqslant i\leqslant 4$};
\draw (49,-23) -- (68,-20);
\draw (49,-23) -- (48,-15.5);
\draw (49,-23) -- (51,-14);
\draw (59.5,-21.5) -- (59,-14);
\draw (68,-20) -- (67,-12.5);
\draw (68,-20) -- (72.5,-13);
\draw (45,-27.5) -- (53,-42);
\draw (46,-38.5) -- (78.5,-34.5);
\draw (70.5,-43.5) -- (75,-24.5);
\node[pt] at (46.5,-30) {};
\node[pt] at (48,-33.5) {};
\node[pt] at (62.5,-36) {};
\node[pt] at (73.5,-29) {};
\node[pt] at (74.5,-25.5) {};
\node[inner sep=1pt] at (92,-23.5) {\scriptsize $\mathbb{P}$};
\end{tikzpicture}
\end{document}
